\section{监控、可观测与治理}

\subsection{部署监控信号}

\begin{table}[H]
\centering
\caption{Deployment 监控清单}
\begin{tabular}{p{0.25\textwidth}p{0.28\textwidth}p{0.25\textwidth}}
\toprule
信号 & 意义 & 触发动作 \\
\midrule
Latency increase & 模型或工具性能退化 & 回退到 last-known-good prompt，补充 optimizer sample \\
Tool error & 工具响应异常 & 暂停相關模块、发出警报、二次验证 \\
Metric drift & 评分下降 & 重新采集训练数据，启动 auto-optimizer \\
User feedback spike & 拒答率/误解率飙升 & 拉入人工审查，更新 evidence trace \\
\bottomrule
\end{tabular}
\end{table}

这些信号通常被 ingested 到 observability stack。Latency 和 Tool error 会让 alerting 立即触发；Metric drift 则需要 schedule 作对比运行（比如 nightly run）。User feedback spike 通常使用 QA note + Slack channel（如 `\#agent-monitors`）做人工 triage。

\begin{warningbox}{治理闭环}
每次重要部署要执行 three-step workflow：define → monitor → annotate。annotate 步骤很关键，确保后续 optimizer 仍记得为什么当初这样配置。
\end{warningbox}

\subsection{Action Timeline}

在每个 release 的 README 中更新 Action Timeline 表，列出模块名称、触发条件、metric threshold 与 responsible owner。如此一来，出现问题时可以快速定位责任链条。

Action Timeline 还可以跟贡献者签名绑定——每次 pipeline 调整要带上 ``who changed what''，便于事故后快速追责。对于多租户团队，建议把每个 timeline entry 也存入 central Incident Log，并用 `Evidence Matrix ID` 做链接。

\subsection{本章小结}

监控与治理是 DSPy pipeline 的第三层护栏，必须把 latency/metric/feedback 等信号转成可执行的响应，而非放任自流。

